% Generated from locale/tr/content/set-theory/replacement/introduction.tex; do not hand-edit.


\olfileid[tr]{sth}{replacement}{intro}
\olsection{Giriş}

Yerine Koyma, $\ZF$ ile~$\Z$ arasındaki farkı yaratan aksiyom şemasıdır.
\crefrange{sth:ordinals::chap}{sth:spine::chap} boyunca onu hiç çekinmeden
kullandık. Bu bölümde sonunda şu soruyu ele alacağız: Yerine Koyma
gerekçelendirilebilir mi?

Soruyu keskinleştirmek için Yerine Koyma'nın gerçekten oldukça
\emph{güçlü} olduğunu gözlemlemekte yarar vardır. Bu bölüm boyunca (ve
\olref[sth][card-arithmetic][fix]{sec} içinde yeniden) ne kadar güçlü
olduğuna ilişkin bir fikir edineceğiz. Fakat bu da gerçekten bir
gerekçelendirme gerektiğini düşündürecektir.

İki tür gerekçelendirmeyi tartışacağız. Kabaca, \emph{dışsal} bir
gerekçelendirme aksiyomu meyveleriyle gerekçelendirme girişimidir;
\emph{içsel} bir gerekçelendirme ise aksiyomun söz konusu matematiksel
kavramlarca doğrulandığını öne sürerek gerekçelendirme girişimidir. Bunun ne
anlama geldiğini bu bölüm boyunca daha iyi kavrayacağız; ancak bu yalnızca
buzdağının görünen kısmıdır. Daha fazlası için özellikle
\citeauthor{Maddy1988a}'ye (\citeyear{Maddy1988a} ve
\citeyear{Maddy1988b}) bakın.

